
\documentclass[11pt]{article}
\usepackage[margin=1in]{geometry}
\usepackage[T1]{fontenc}
\usepackage[utf8]{inputenc}
\usepackage{lmodern}
\usepackage{amsmath}
\usepackage{amssymb}
\usepackage{graphicx}
\usepackage{booktabs}
\usepackage{float}
\usepackage{hyperref}
\hypersetup{hidelinks}
\setlength{\parskip}{0.5em}
\setlength{\parindent}{0pt}
\begin{document}
\sloppy
\section{Step 6: Trial-Level ROI Analysis Adjusted for Question Covariates}
\label{sec:step6_covariates}

\subsection{Objective}
The objective of Step~6 is to test whether the abstract-versus-concrete effect in the onset-locked
left frontal ROI remains after accounting for differences between:
\begin{itemize}
    \item participants;
    \item questions;
    \item question length and question structure.
\end{itemize}

This step is motivated by two observations from the previous analyses:
\begin{enumerate}
    \item the Step~5 fixed-window front ROI HbO result was the strongest inferential signal so far,
    but it did not reach conventional statistical significance;
    \item the onset-locked plots suggest meaningful heterogeneity across individual trials,
    across channels, and across local channel waveforms.
\end{enumerate}

Therefore, Step~6 should ask:
\begin{quote}
after controlling for question size and structure, and after accounting for participant-level and
question-level variability, is there evidence of an abstract-versus-concrete effect in the left
frontal ROI?
\end{quote}

\subsection{Why This Step Is Needed}
In the current dataset, trials differ not only by condition but also by the properties of the question itself.
For example, questions vary in:
\begin{itemize}
    \item total number of words;
    \item number of sentences;
    \item average sentence length.
\end{itemize}

These properties may influence reading load, linguistic complexity, and cognitive effort.
If abstract questions are systematically longer or more linguistically dense than concrete questions,
part of the observed fNIRS difference may reflect question size or structure rather than condition alone.

At the same time, some participants may generally show larger or smaller fNIRS responses than others,
and some questions may generally evoke larger or smaller responses than others.
A trial-level model that accounts for both participants and questions is therefore the most direct
next step.

\subsection{Core Analytical Decision}
To keep Step~6 disciplined and interpretable, the following should remain fixed from Step~5:
\begin{enumerate}
    \item the same Step~5 included participant-session cohort;
    \item the same preprocessed fNIRS files from Step~3;
    \item the same long-channel-only rule;
    \item the same front ROI definition;
    \item the same primary onset-locked epoching scheme;
    \item the same primary fixed HbO window:
    \begin{verbatim}
    7.0 s to 11.0 s after question onset
    \end{verbatim}
\end{enumerate}

Thus, Step~6 changes only one main thing:
\begin{quote}
the analysis moves from unadjusted condition comparisons to adjusted trial-level modelling
with question covariates and random effects.
\end{quote}

\subsection{Primary Scientific Question}
The primary Step~6 question is:
\begin{quote}
In the left frontal ROI, does the abstract-versus-concrete effect remain after adjusting for
question length and participant/question variability?
\end{quote}

\subsection{Primary Outcome}
The primary outcome remains the same as in Step~5:
\begin{quote}
the trial-level left frontal ROI HbO mean in the fixed 7--11~s onset-locked window.
\end{quote}

This is kept fixed so that Step~6 isolates the impact of the covariate adjustment rather than
changing both the dependent variable and the model at the same time.

\subsection{Input Variables}
Step~6 combines information from:
\begin{itemize}
    \item Step~3 preprocessed fNIRS files;
    \item Step~3 harmonized trigger table;
    \item Step~5 onset-locked trial table;
    \item Step~2 question metadata.
\end{itemize}

At minimum, the question-level covariates available for this step are:
\begin{itemize}
    \item \texttt{total\_word\_count};
    \item \texttt{sentence\_count};
    \item \texttt{sentence\_length};
    \item \texttt{correctness} from the original ENEM item-level metadata.
\end{itemize}

\subsection{Important Collinearity Rule}
The variables \texttt{total\_word\_count}, \texttt{sentence\_count}, and \texttt{sentence\_length}
are structurally related.
In particular, average sentence length is approximately:
\begin{equation}
\texttt{sentence\_length}
\approx
\frac{\texttt{total\_word\_count}}{\texttt{sentence\_count}}.
\end{equation}

Therefore, Step~6 should \textbf{not} include all three structure variables in the same model.
Instead, the analysis should use:
\begin{itemize}
    \item one \textbf{primary adjusted model} with \texttt{total\_word\_count};
    \item two \textbf{alternative adjusted models}, one with \texttt{sentence\_count} and one with
    \texttt{sentence\_length}.
\end{itemize}

This keeps the covariate adjustment interpretable and avoids unnecessary multicollinearity.

\subsection{Primary Adjustment Strategy}
The primary covariate for Step~6 should be:
\begin{quote}
\textbf{log-transformed total word count}
\end{quote}

This is the best first control because it is the most direct measure of overall reading load and
is likely to be skewed in raw form.

Thus, the primary covariate transformation should be:
\begin{equation}
w_k = z\!\left(\log(\texttt{total\_word\_count}_k)\right),
\end{equation}
where \(z(\cdot)\) denotes standardization to mean zero and standard deviation one.

\subsection{Modeling Strategy}
The basic unit of analysis in Step~6 is:
\begin{quote}
one participant-question trial.
\end{quote}

Each trial belongs simultaneously to:
\begin{itemize}
    \item one participant;
    \item one question.
\end{itemize}

For this reason, the natural model for Step~6 is a \textbf{trial-level mixed-effects model}
with random intercepts for:
\begin{itemize}
    \item participant;
    \item question.
\end{itemize}

\subsection{Why a Mixed-Effects Model Is Appropriate}
A mixed-effects model is appropriate because it directly addresses the two main sources of nuisance variation:
\begin{itemize}
    \item some participants may systematically show larger or smaller responses;
    \item some questions may systematically evoke larger or smaller responses.
\end{itemize}

By including both participant and question random intercepts, the model can estimate the condition
effect after accounting for these differences.

\subsection{Primary Step~6 Model}
Let \(y_{ik}\) denote the trial-level front ROI HbO value for participant \(i\) on question \(k\)
in the fixed 7--11~s onset-locked window.
Let \(A_{ik}\) be the condition indicator:
\begin{equation}
A_{ik} =
\begin{cases}
1 & \text{if the trial is Abstract} \\
0 & \text{if the trial is Concrete}
\end{cases}
\end{equation}

Let \(w_k\) be the standardized log word-count covariate.

The primary Step~6 model should be:
\begin{equation}
y_{ik}
=
\beta_0
+
\beta_1 A_{ik}
+
\beta_2 w_k
+
u_i
+
v_k
+
\varepsilon_{ik},
\end{equation}
where:
\begin{itemize}
    \item \(u_i \sim \mathcal{N}(0, \sigma_u^2)\) is the participant random intercept;
    \item \(v_k \sim \mathcal{N}(0, \sigma_v^2)\) is the question random intercept;
    \item \(\varepsilon_{ik} \sim \mathcal{N}(0, \sigma^2)\) is the residual term.
\end{itemize}

\subsection{Primary Inferential Target}
The primary inferential target in Step~6 is:
\begin{equation}
\beta_1
\end{equation}
which represents the adjusted abstract-versus-concrete effect in the front ROI after controlling
for word count and participant/question variability.

\subsection{Primary Hypothesis}
The primary hypothesis is:

\paragraph{Null hypothesis}
\begin{equation}
H_0: \beta_1 = 0
\end{equation}

\paragraph{Alternative hypothesis}
\begin{equation}
H_1: \beta_1 \neq 0
\end{equation}

The significance threshold should remain:
\begin{verbatim}
alpha = 0.05
\end{verbatim}

\subsection{Alternative Adjusted Models}
To assess whether the result depends on the specific structural covariate chosen, Step~6 should also fit:

\subsubsection{Model A: Sentence-count model}
\begin{equation}
y_{ik}
=
\beta_0
+
\beta_1 A_{ik}
+
\beta_2 s_k
+
u_i
+
v_k
+
\varepsilon_{ik},
\end{equation}
where:
\begin{equation}
s_k = z(\texttt{sentence\_count}_k)
\end{equation}

\subsubsection{Model B: Sentence-length model}
\begin{equation}
y_{ik}
=
\beta_0
+
\beta_1 A_{ik}
+
\beta_2 \ell_k
+
u_i
+
v_k
+
\varepsilon_{ik},
\end{equation}
where:
\begin{equation}
\ell_k = z(\texttt{sentence\_length}_k)
\end{equation}

These alternative models are not replacements for the primary word-count model.
They are robustness checks.

\subsection{Difficulty-Adjusted Sensitivity Model}
Because Step~2 also contains the ENEM item-level correctness value, Step~6 should include one
additional sensitivity model that adjusts for both question size and question difficulty.

Let:
\begin{equation}
c_k = z(\texttt{correctness}_k)
\end{equation}

Then the difficulty-adjusted sensitivity model should be:
\begin{equation}
y_{ik}
=
\beta_0
+
\beta_1 A_{ik}
+
\beta_2 w_k
+
\beta_3 c_k
+
u_i
+
v_k
+
\varepsilon_{ik}.
\end{equation}

This model is useful because it tests whether the abstract-versus-concrete effect remains after
adjusting both for question length and for the original population-level item correctness.

\subsection{Variables That Should Not Be in the Primary Model}
To keep the Step~6 model interpretable, the following variables should \textbf{not} be included
in the primary model:
\begin{itemize}
    \item response time;
    \item participant correctness on the trial;
    \item all three question-structure variables at once.
\end{itemize}

The reason is that response time and participant correctness may be \textbf{mediators} rather than pure confounders.
For example, if abstract questions truly require more processing, then controlling for response time
could remove part of the effect of interest.

\subsection{Optional Sensitivity Models Involving Behavior}
Response time and participant correctness may still be useful in \textbf{sensitivity analyses}, but they
must be explicitly labeled as such.

\subsubsection{Sensitivity model with response time}
Let \(r_{ik}\) be standardized log response time.
Then:
\begin{equation}
y_{ik}
=
\beta_0
+
\beta_1 A_{ik}
+
\beta_2 w_k
+
\beta_3 r_{ik}
+
u_i
+
v_k
+
\varepsilon_{ik}
\end{equation}

\subsubsection{Sensitivity model with participant correctness}
Let \(p_{ik}\) be the participant correctness indicator.
Then:
\begin{equation}
y_{ik}
=
\beta_0
+
\beta_1 A_{ik}
+
\beta_2 w_k
+
\beta_3 p_{ik}
+
u_i
+
v_k
+
\varepsilon_{ik}
\end{equation}

These should be interpreted carefully because they may adjust away part of the condition effect.

\subsection{Balance Check Before Modeling}
Before fitting the mixed-effects models, Step~6 should include a simple item-level balance check
between abstract and concrete questions.

The balance check should compare the two condition groups on:
\begin{itemize}
    \item \texttt{total\_word\_count};
    \item \texttt{sentence\_count};
    \item \texttt{sentence\_length};
    \item \texttt{correctness}.
\end{itemize}

This balance check should be primarily descriptive and should report:
\begin{itemize}
    \item mean;
    \item standard deviation;
    \item median;
    \item minimum;
    \item maximum;
    \item standardized mean difference.
\end{itemize}

Optional simple two-group tests may also be reported, but the emphasis should remain descriptive.

\subsection{Step-by-Step Workflow}

\subsubsection{Step~6.1: Freeze the Step~5 cohort}
Reuse the same included participant-session set as Step~5.
Do not redefine the cohort unless a documented data problem requires it.

\subsubsection{Step~6.2: Build the trial-level analysis table}
Create one row per valid participant-question trial used in Step~5.
The trial table should include at least:
\begin{itemize}
    \item \texttt{participant\_id};
    \item \texttt{question\_id};
    \item \texttt{condition} (\texttt{abstract} or \texttt{concrete});
    \item front ROI HbO 7--11~s mean;
    \item back ROI HbO 7--11~s mean;
    \item \texttt{total\_word\_count};
    \item \texttt{sentence\_count};
    \item \texttt{sentence\_length};
    \item \texttt{correctness} from ENEM;
    \item response time;
    \item participant correctness.
\end{itemize}

\subsubsection{Step~6.3: Run the item-level balance check}
Produce a descriptive summary comparing abstract and concrete questions on the item covariates.

\subsubsection{Step~6.4: Prepare the covariates}
Apply the following transformations:
\begin{itemize}
    \item \texttt{log(total\_word\_count)};
    \item standardize continuous covariates;
    \item code condition as binary;
    \item verify no impossible or missing values remain in the primary model table.
\end{itemize}

\subsubsection{Step~6.5: Fit the primary word-count adjusted mixed model}
Fit the front ROI HbO model with:
\begin{itemize}
    \item fixed effect of condition;
    \item fixed effect of log word count;
    \item random intercept for participant;
    \item random intercept for question.
\end{itemize}

\subsubsection{Step~6.6: Fit the sentence-count and sentence-length alternative models}
Fit the two alternative front ROI models:
\begin{itemize}
    \item condition + sentence count;
    \item condition + sentence length.
\end{itemize}

\subsubsection{Step~6.7: Fit the difficulty-adjusted sensitivity model}
Fit the front ROI model including:
\begin{itemize}
    \item condition;
    \item log word count;
    \item ENEM correctness.
\end{itemize}

\subsubsection{Step~6.8: Fit the back ROI model}
Repeat the primary word-count adjusted model using the back ROI HbO response instead of the front ROI response.
This is a secondary model and is useful for checking whether any adjusted effect remains primarily frontal.

\subsubsection{Step~6.9: Fit behavioral sensitivity models}
Optionally fit:
\begin{itemize}
    \item condition + word count + response time;
    \item condition + word count + participant correctness.
\end{itemize}

These should be labeled as sensitivity analyses only.

\subsubsection{Step~6.10: Summarize and compare the model results}
Collect the following for each model:
\begin{itemize}
    \item estimate of the condition coefficient \(\beta_1\);
    \item standard error;
    \item test statistic;
    \item p-value;
    \item 95\% confidence interval;
    \item residual variance;
    \item participant random-intercept variance;
    \item question random-intercept variance.
\end{itemize}

\subsection{Primary Reporting Quantities}
The primary Step~6 results table must report:
\begin{itemize}
    \item number of included participants;
    \item number of included questions;
    \item number of included trials;
    \item front ROI adjusted condition coefficient;
    \item standard error;
    \item test statistic;
    \item p-value;
    \item 95\% confidence interval;
    \item word-count coefficient;
    \item participant random-intercept variance;
    \item question random-intercept variance.
\end{itemize}

\subsection{Secondary Reporting Quantities}
The secondary Step~6 results tables should report:
\begin{itemize}
    \item sentence-count model results;
    \item sentence-length model results;
    \item difficulty-adjusted model results;
    \item back ROI model results;
    \item behavioral sensitivity model results.
\end{itemize}

\subsection{Primary Conclusion Rule}
The final Step~6 conclusion should be based first on the condition coefficient in the
primary word-count adjusted front ROI model.

The interpretation rule should be:
\begin{itemize}
    \item if the adjusted condition coefficient is significant and positive, conclude that abstract
    questions are associated with a larger front ROI HbO response than concrete questions even after
    adjusting for question length and participant/question variability;
    \item if the adjusted condition coefficient is significant and negative, conclude that concrete
    questions are associated with a larger front ROI HbO response than abstract questions after adjustment;
    \item if the adjusted condition coefficient is not significant, conclude that the adjusted model
    does not show statistically significant evidence for a front ROI abstract-versus-concrete effect.
\end{itemize}

\subsection{Interpretation Templates}

\paragraph{Template if the adjusted effect is significant and abstract \(>\) concrete.}
In the trial-level mixed-effects analysis adjusted for log word count, abstract questions were
associated with a significantly larger left frontal ROI HbO response than concrete questions
(\(\beta=\ldots\), \(SE=\ldots\), \(p=\ldots\), 95\% CI \([\ldots,\ldots]\)).
This indicates that the frontal condition effect remains after accounting for both participant-level
and question-level variability and for differences in question length.

\paragraph{Template if the adjusted effect is significant and concrete \(>\) abstract.}
In the trial-level mixed-effects analysis adjusted for log word count, concrete questions were
associated with a significantly larger left frontal ROI HbO response than abstract questions
(\(\beta=\ldots\), \(SE=\ldots\), \(p=\ldots\), 95\% CI \([\ldots,\ldots]\)).

\paragraph{Template if the adjusted effect is not significant.}
In the trial-level mixed-effects analysis adjusted for log word count, the abstract-versus-concrete
effect in the left frontal ROI was not statistically significant
(\(\beta=\ldots\), \(SE=\ldots\), \(p=\ldots\), 95\% CI \([\ldots,\ldots]\)).
This suggests that, after accounting for participant differences, question differences, and question
length, the current data do not provide statistically significant evidence for a frontal condition effect.

\subsection{Expected Outputs}

\subsubsection*{Tables}
\begin{verbatim}
01_step6_item_covariate_table.csv
02_step6_item_balance_summary.csv
03_step6_trial_model_table.csv
04_step6_primary_front_wordcount_model.csv
05_step6_front_sentencecount_model.csv
06_step6_front_sentencelength_model.csv
07_step6_front_wordcount_difficulty_model.csv
08_step6_back_wordcount_model.csv
09_step6_behavior_sensitivity_models.csv
10_step6_exclusion_log.csv
11_step6_final_conclusion.csv
\end{verbatim}

\subsubsection*{Figures}
\begin{verbatim}
figures/step6/wordcount_by_condition.png
figures/step6/sentencecount_by_condition.png
figures/step6/sentencelength_by_condition.png
figures/step6/enem_correctness_by_condition.png
figures/step6/front_roi_vs_wordcount_scatter.png
figures/step6/front_roi_partial_effect_condition.png
figures/step6/model_residuals_primary.png
figures/step6/model_random_intercepts.png
\end{verbatim}

\subsubsection*{Reports}
\begin{verbatim}
reports/step6/step6_covariate_adjusted_report.tex
reports/step6/tables/item_balance_table.tex
reports/step6/tables/primary_model_table.tex
reports/step6/tables/secondary_model_table.tex
reports/step6/text/final_conclusion.tex
\end{verbatim}

\subsection{Minimal Figures for the Main Report}
The main Step~6 report should include at least:
\begin{enumerate}
    \item a balance plot for word count by condition;
    \item a balance plot for sentence count by condition;
    \item a balance plot for sentence length by condition;
    \item a scatter plot of front ROI HbO versus word count, colored by condition;
    \item a partial-effect plot showing the adjusted condition effect from the primary model;
    \item one simple residual diagnostic plot for the primary model.
\end{enumerate}

\subsection{Quality-Control Checks}
Before Step~6 is considered complete, verify that:
\begin{enumerate}
    \item the Step~5 cohort is unchanged unless explicitly re-reviewed;
    \item the front ROI 7--11~s HbO outcome is identical to the Step~5 trial-level outcome;
    \item condition labels are complete and valid;
    \item the question covariates are correctly merged by question ID;
    \item no impossible word-count or sentence-count values remain;
    \item the primary model contains only one structural question covariate;
    \item sentence length is not used in the same model as both word count and sentence count;
    \item the participant and question random effects are identifiable;
    \item the primary word-count adjusted model converges;
    \item all alternative and sensitivity models are clearly labeled as non-primary.
\end{enumerate}

\subsection{Fallback Rule if the Mixed Model Fails}
If the primary mixed-effects model fails to converge reliably, the fallback analysis should be:
\begin{quote}
a linear model with participant fixed effects and question fixed effects, keeping the same
front ROI outcome and the same primary covariate.
\end{quote}

This fallback should be used only if needed and must be documented explicitly.

\subsection{Minimal Completion Criteria}
Step~6 is complete when:
\begin{itemize}
    \item the item-level balance check has been completed;
    \item the trial-level analysis table has been built;
    \item the primary word-count adjusted front ROI model has been fit;
    \item the alternative structural-covariate models have been fit;
    \item the difficulty-adjusted sensitivity model has been fit;
    \item the back ROI model has been fit;
    \item the primary adjusted p-value has been reported;
    \item a short final conclusion has been generated;
    \item all exclusions and warnings have been documented.
\end{itemize}

\subsection{Short Practical Summary}
In simple terms, Step~6 should:
\begin{quote}
keep the same front ROI HbO 7--11~s outcome from Step~5, move to a trial-level mixed-effects model,
control first for total word count, account for different participants and different questions,
then check whether the abstract-versus-concrete effect survives that adjustment.
\end{quote}
\end{document}
